\documentclass[11pt]{article}
\usepackage{amsmath,amssymb,amsthm,mathtools}
\usepackage[margin=1in]{geometry}
\usepackage{booktabs}
\usepackage{hyperref}
\usepackage{enumitem}
\hypersetup{pdftitle={The Thomson problem for N=14: a computer-assisted proof},pdfauthor={}}

\newtheorem{theorem}{Theorem}[section]
\newtheorem{proposition}[theorem]{Proposition}
\newtheorem{lemma}[theorem]{Lemma}
\newtheorem{corollary}[theorem]{Corollary}
\theoremstyle{definition}
\newtheorem{definition}[theorem]{Definition}
\newtheorem{remark}[theorem]{Remark}

\newcommand{\R}{\mathbb{R}}
\newcommand{\Q}{\mathbb{Q}}
\newcommand{\Sph}{\mathbb{S}^2}
\newcommand{\ip}[2]{\langle #1,#2\rangle}
\newcommand{\Xs}{X^{\ast}}
\newcommand{\Es}{E^{\ast}}
\newcommand{\hs}{h^{\ast}}
\newcommand{\PSD}{\succeq 0}
% ---- numbers of the Case III certificate (filled in from the verification logs)
\newcommand{\TUBE}{1.95\cdot10^{-4}}        % verified tube radius tau (1.9482e-4 rounded up)
\newcommand{\THLOC}{4\cdot10^{-14}}
\newcommand{\THFAR}{10^{-12}}
\newcommand{\DELTAZ}{10^{-15}}
\newcommand{\CMAX}{237.3}          % largest c_M
\newcommand{\NPPR}{533\,876}\newcommand{\NPRR}{630\,673}\newcommand{\NRRR}{475\,861}

\title{The Thomson problem for $N=14$:\\ the bicapped hexagonal antiprism is optimal\\[0.3em]
\large a computer-assisted proof}
\author{}
\date{Public verification package, 1 October 2026}

\begin{document}
\maketitle

\begin{abstract}
Let $\Xs$ be the bicapped hexagonal antiprism with $D_{6d}$ symmetry. We prove that $\Xs$ is, up to isometries, the unique
minimizer of the Coulomb energy $E(X)=\sum_{i<j}|x_i-x_j|^{-1}$ among $14$-point configurations on $\Sph$;
$\Es=E(\Xs)=69.30636329662642125848999127638\ldots$ (rigorous enclosure).
The proof is computer-assisted. Three-point semidefinite bounds, converted into exact rational certificates and verified by
branch and bound in outward-rounded interval arithmetic, show: (0) a competitor has no two points with inner product
$\ge0.8$; (I), (II) a competitor contains a pair of points with inner product $<-0.99$; (III) a certificate with an
energy offset just below $\Es$ forces all inner products of such a competitor to lie within $\TUBE$ of
the corresponding inner products of $\Xs$. A rigidity lemma and a quantitative second-order analysis at $\Xs$ complete the
proof.
Small residual values and gradients at the candidate triple types are included in the interval estimates.
A counting argument, using the twelve triples through each pair, supplies the quantitative localization bound.
\end{abstract}


%======================================================================
\section{Introduction and main result}
%======================================================================

For a finite configuration $X=\{x_1,\dots,x_N\}\subset\Sph$ of distinct points let
\[
  E(X)=\sum_{1\le i<j\le N}\frac{1}{|x_i-x_j|}=\sum_{i<j}\varphi(\ip{x_i}{x_j}),\qquad \varphi(t)=(2-2t)^{-1/2}.
\]
Thomson's problem asks for the minimizers of $E$ among $N$-point configurations. We treat $N=14$ by combining
three-point positive semidefinite kernels with interval estimates and a geometric rigidity argument.
The kernel method is related to the spherical-code bounds of Bachoc and Vallentin~\cite{BV} and the energy bounds of
Cohn and Woo~\cite{CohnWoo}. We give the positivity and counting identities needed here explicitly.
The numerical optimization used to find the certificates is not an assumption of the proof: the resulting rational
data are checked separately.

\paragraph{Numerical conventions.}
Terminating decimals denote exact rationals, except for the stored case cutoffs $-0.99$, $-0.98$, $0.8$
and slack thresholds $0.0459$, $4\cdot10^{-14}$ and $10^{-12}$, which denote their nearest binary64 values.
The same stored cutoffs are used on both sides of every case distinction. In particular,
$0.8$ denotes $3602879701896397/4503599627370496$ in these case conditions.
Approximate decimal values are marked by $\approx$ or an ellipsis.

\paragraph{The candidate.}
For $h\in(0,1)$ let $r=\sqrt{1-h^2}$ and let $X(h)$ consist of the two poles $\pm e_3$, the six points
$(r\cos\frac{k\pi}{3},r\sin\frac{k\pi}{3},h)$ ($k=0,\dots,5$), and the six points
$(r\cos(\frac{k\pi}{3}+\frac{\pi}{6}),r\sin(\frac{k\pi}{3}+\frac{\pi}{6}),-h)$ ($k=0,\dots,5$): two parallel regular hexagons,
rotated against each other by $30^\circ$, capped by the poles. Its energy is
\begin{equation}\label{eq:Efam}
 E(X(h))=\tfrac12+\frac{12}{\sqrt{2-2h}}+\frac{12}{\sqrt{2+2h}}+\frac{15+12/\sqrt3}{r}
 +\sum_{c\in\{\frac{\sqrt3}{2},0,-\frac{\sqrt3}{2}\}}\frac{12}{\sqrt{2r^2(1-c)+4h^2}} .
\end{equation}

\begin{lemma}[the candidate]\label{lem:hstar}
The function $h\mapsto E(X(h))$ has a critical point $\hs$ with
\begin{align*}
 \hs &\in 0.45536779516246300350924502372107445741075259468695 \pm 2\cdot 10^{-51},\\
 \Es:=E(X(\hs)) &\in 69.306363296626421258489991276382007864282129573759 \pm 10^{-49}.
\end{align*}
Moreover $\Xs:=X(\hs)$ is a critical point of $E$ on $(\Sph)^{14}$.
\end{lemma}
\begin{proof}
Let $h_0$ be the terminating decimal displayed in the enclosure of $\hs$, and put
$I_0=[h_0-10^{-40},h_0+10^{-40}]$. The program \path{Certificates/common/enclose_h.py}, using
\texttt{Arb} at 256 bits, verifies that the derivative of~\eqref{eq:Efam} is negative at the left endpoint,
positive at the right endpoint, and has positive derivative throughout $I_0$. Thus a unique critical height exists in
$I_0$. The interval Newton image
$h_0-E'(h_0)/E''(I_0)$ is contained in $I_0$ and encloses this height.
Direct interval evaluation on that image verifies both displayed enclosures.
Criticality of $\Xs$: the configuration $\Xs$ is invariant under the group $D_{6d}$, which acts transitively on each of the
orbits $\{\pm e_3\}$ and on the twelve ring points.  The gradient $\nabla_{x_i}E(\Xs)$ (projected to $T_{x_i}\Sph$) is fixed by
the stabiliser of $x_i$.  For a pole the stabiliser contains the rotation by $60^\circ$ about $e_3$, whose only fixed tangent
vector is $0$.  For a ring point the stabiliser contains the reflection in the vertical plane through $x_i$, so the tangential
gradient is a multiple $g$ of the meridian direction, and $g$ is the same for all twelve ring points up to the symmetry
(the improper rotation $S_{12}$ exchanges the rings and reverses the vertical direction).  Consequently the derivative of
$E$ along the family $X(h)$ equals $12\,g\cdot\frac{d}{dh}(\text{arc length})$, i.e.\ $g$ is a nonzero multiple of
$\frac{d}{dh}E(X(h))$, which vanishes at $\hs$.  Hence $\nabla E(\Xs)=0$.
\end{proof}

\begin{theorem}[optimality]\label{thm:main}
For every configuration $X$ of $14$ distinct points on $\Sph$ one has $E(X)\ge\Es$, with equality if and only if
$X=Q\Xs$ for some $Q\in O(3)$.
\end{theorem}

The proof combines the following statements, which are of independent interest.

\begin{theorem}[configurations without a nearly antipodal pair]\label{thm:A}
Every configuration $X$ of $14$ distinct points on $\Sph$ with $m(X)=\min_{i\ne j}\ip{x_i}{x_j}\ge-0.99$ satisfies
$E(X)\ge\Es+3.6\cdot10^{-4}$. Every configuration containing two points with $\ip{x_i}{x_j}\ge0.8$ satisfies
$E(X)\ge\Es+0.086$.
\end{theorem}

\begin{theorem}[local optimality]\label{thm:B}
$\Xs$ is a strict local minimizer of $E$: in the gauge coordinates $p$ of Section~\ref{sec:local},
$E(p)>\Es$ for $0<|p-p^\ast|_2\le0.0075$. Moreover (Lemma~\ref{lem:rigidity}) every configuration whose inner products
are within $\tau\le10^{-3}$ of those of $\Xs$ (in the sense made precise there) satisfies $|p-p^\ast|_2\le23\,\tau$.
\end{theorem}

\begin{theorem}[the cap]\label{thm:C}
Let $X$ be a configuration with $m(X)\le-0.99$, all of whose inner products are smaller than $0.8$, and $E(X)\le\Es$. Then
$X$ satisfies the hypotheses (i)--(iii) of Lemma~\ref{lem:rigidity} with $\tau=\TUBE$; in particular every inner product
of $X$ lies within $\TUBE$ of an inner product of the corresponding pair class of $\Xs$.
\end{theorem}

The nine distinct inner products occurring in $\Xs$ are
\begin{equation}\label{eq:nodes}
 -1,\ \pm\hs,\ \tfrac{1+h^{2}}{2},\ \tfrac{3h^2-1}{2},\ 2h^2-1,\ -h^2+\tfrac{\sqrt3}{2}(1-h^2),\ -h^2,\ -h^2-\tfrac{\sqrt3}{2}(1-h^2)
 \qquad(h=\hs),
\end{equation}
numerically $-1,\ \pm0.455368,\ 0.603680,\ -0.188960,\ -0.585280,\ 0.479087,\ -0.207360,\ -0.893806$.
In particular $m(\Xs)=-1$ is attained exactly by the pole pair, and every other inner product of $\Xs$ is $\ge -0.8939$.

%======================================================================
\section{Strategy of the proof}\label{sec:strategy}
%======================================================================
By compactness of $(\Sph)^{14}$ and since $E\to\infty$ when two points collide, a minimizer exists. Let $X$ be any
configuration with $E(X)\le\Es$, let $m=m(X)$ be its minimal inner product, attained by a pair $\{P,Q\}$ (the
\emph{minimal pair}; we fix one if several pairs attain $m$), and let $M(X)$ be its maximal inner product.
\begin{enumerate}[label=(\Roman*)]
\item[(0)] $M(X)\ge0.8$ is impossible: an untyped three-point bound together with a counting argument gives
      $E(X)\ge\Es+0.086$ (Proposition~\ref{prop:sep}).
\item $m\ge-0.98$ is impossible: an untyped three-point bound (Section~\ref{sec:case1}) gives $E(X)\ge e_{\rm I}>\Es$.
\item $-0.99\le m\le-0.98$ is impossible: a typed three-point bound with shares ($P/R$ typing,
      Section~\ref{sec:slabs}) gives $E(X)\ge e_{\rm II}>\Es$.
\item $-1\le m\le-0.99$ (the \emph{cap}) and $M(X)<0.8$. This class contains $\Xs$.
      A typed certificate ($P/R$ typing, degree $20$), with an energy offset just below $\Es$ and a quantitative
      positive lower bound away from the triple types of $\Xs$, shows that
      all inner products of $X$ are within $\TUBE$ of those of $\Xs$ (Theorem~\ref{thm:C}, Section~\ref{sec:cap}).
      The rigidity lemma (Section~\ref{sec:rigidity}) then puts $X$, after an orthogonal change of coordinates, into a
      neighbourhood of $\Xs$ in which $\Xs$ is the unique minimizer by a quantitative second-order argument
      (Section~\ref{sec:local}).
\end{enumerate}
\begin{proof}[Proof of Theorem~\ref{thm:main}]
Let $E(X)\le\Es$. By Theorem~\ref{thm:A}, $M(X)<0.8$ and $m(X)<-0.99$. By Theorem~\ref{thm:C} the hypotheses of
Lemma~\ref{lem:rigidity} hold with $\tau=\TUBE$, hence $|p-p^\ast|_2\le23\tau=0.004485<0.0075$ after an
orthogonal map, and Proposition~\ref{prop:local} gives $p=p^\ast$, i.e.\ $X=Q\Xs$ with $Q\in O(3)$, and in particular
$E(X)=\Es$. Hence no configuration has energy below $\Es$, and the configurations of energy $\Es$ are exactly the images
of $\Xs$ under $O(3)$.
\end{proof}

%======================================================================
\section{Three-point bounds with typed kernels}\label{sec:bounds}
%======================================================================

\subsection{Positive definite kernels relative to a root point}
Write $u=\ip{x}{y}$, $v=\ip{x}{z}$, $t=\ip{y}{z}$ for three points $x,y,z\in\Sph$ and define polynomials
$Q_k(u,v,t)$ by $Q_0=1$, $Q_1=t-uv$ and
\[
 Q_k=2(t-uv)\,Q_{k-1}-(1-u^2)(1-v^2)\,Q_{k-2}\qquad(k\ge2).
\]
If $|u|,|v|<1$ then $Q_k(u,v,t)=\big((1-u^2)(1-v^2)\big)^{k/2}\cos(k\theta)$, where $\theta$ is the angle between the
projections of $y$ and $z$ to the plane $x^\perp$. For $k,i,j\ge0$ let $T_i$ denote the Chebyshev polynomials.

\begin{lemma}[positivity]\label{lem:pd}
Let $x\in\Sph$, let $Y\subset\Sph$ be finite, let $\tau:Y\to\{1,\dots,T\}$ be arbitrary ("types"), for every type
$\sigma$ and every $k$ let $p^{(k)}_{\sigma,0},\dots,p^{(k)}_{\sigma,n_{k,\sigma}-1}$ be real polynomials, and let
$F^{(k)}=\big(F^{(k)}_{(\sigma,i),(\sigma',j)}\big)$ be real symmetric positive semidefinite matrices. Then
\[
 \sum_{y,z\in Y}\ \sum_{k\ge 0}\ \sum_{i,j} F^{(k)}_{(\tau(y),i),(\tau(z),j)}\,p^{(k)}_{\tau(y),i}(\ip{x}{y})\,p^{(k)}_{\tau(z),j}(\ip{x}{z})\,
 Q_k(\ip{x}{y},\ip{x}{z},\ip{y}{z})\ \ge\ 0 ,
\]
where the sum runs over all ordered pairs, including $y=z$ and $y=x$ or $z=x$ if $x\in Y$.
\end{lemma}
\begin{proof}
For $k=0$ put $w_y=1$ for every $y$. For $k\ge1$ and $y\neq\pm x$ write
$w_y=(1-\ip{x}{y}^2)^{k/2}e^{\mathrm{i}k\alpha_y}$ with $\alpha_y$ the azimuth of $y$ around $x$,
and $w_y=0$ if $y=\pm x$ and $k\ge1$. Then $Q_k(\ip{x}{y},\ip{x}{z},\ip{y}{z})=\operatorname{Re}(w_y\overline{w_z})$ for all
$y,z$ (by continuity also in the degenerate cases), and the inner sum equals
$\operatorname{Re}\,\xi^{*}F^{(k)}\xi\ge 0$ with $\xi=\sum_{y}w_y\,e_{\tau(y)}\otimes\big(p^{(k)}_{\tau(y),i}(\ip{x}{y})\big)_i$.
\end{proof}
Our certificates use $p^{(k)}_{\sigma,i}=T_i$, with $0\le i\le K-k$.
In the cap certificate only, a kernel rooted at a pole uses the smaller pole-type basis
$0\le i<\min(K-k+1,2)$; its arguments of that type are evaluated only at $\ip PQ$ and $1$.
All other type blocks use the full basis.

\subsection{Types, pair classes, shares and slacks}
Fix a partition of the $14$ points into types $1,\dots,T$ with prescribed cardinalities $n_1,\dots,n_T$. For a root type
$X$ and $k=0,\dots,K$ let $F_X^{(k)}\PSD$ be as in Lemma~\ref{lem:pd} ($n_k=K-k+1$), and for a point $x$ of type $X$ and
points $y,z$ define the kernel
\[
 K_x(y,z)=\sum_{k=0}^{K}\sum_{i,j}F^{(k)}_{X,(\tau(y),i),(\tau(z),j)}\,T_i(u)\,T_j(v)\,Q_k(u,v,t).
\]
For a root type $X$ and a type $\sigma$ put
\begin{align*}
 g_X(s;\sigma)&=\sum_{k}\ \big(T(s)^{\!\top}F^{(k)}_{X,\sigma\sigma}T(s)\big)(1-s^2)^k
 +2\sum_{i,j}F^{(0)}_{X,(X,i),(\sigma,j)}T_i(1)T_j(s),\\
 \mathrm{sing}_X&=\sum_{i,j}F^{(0)}_{X,(X,i),(X,j)}
\end{align*}
(the contributions of the degenerate pairs $y=z$ and $\{y,z\}\ni x$), and for a pair class $\{a,b\}$ of types
\[
 \Psi_{ab}(s)=\varphi(s)-g_a(s;b)-g_b(s;a).
\]
For every pair class $\{a,b\}$ and every type $c$ we choose a univariate function $\psi_{ab\to c}$ (the \emph{share} that a
pair of types $a,b$ gives to a third vertex of type $c$) such that
\begin{equation}\label{eq:sharesum}
 \sum_c m_c(ab)\,\psi_{ab\to c}=\Psi_{ab},\qquad m_c(ab)=n_c-[c=a]-[c=b]
\end{equation}
($m_c(ab)$ is the number of third vertices of type $c$ available to a pair of types $a,b$). In our certificates all shares but
one per class are polynomials, and the remaining one is determined by \eqref{eq:sharesum}; it carries the singular part of
$\varphi$. Finally we choose constants $c_{abc}$ for every pattern $\{a,b,c\}$ of three types with
\begin{equation}\label{eq:constsum}
 \sum_{\text{triples }T}c_{T}=e+\sum_{x}\mathrm{sing}_{\tau(x)} .
\end{equation}
For a triple $T=\{x,y,z\}$ of points define the \emph{slack}
\begin{equation}\label{eq:slack}
\begin{aligned}
 \lambda_T={}&\psi_{\tau(y)\tau(z)\to\tau(x)}(\ip{y}{z})+\psi_{\tau(x)\tau(z)\to\tau(y)}(\ip{x}{z})
 +\psi_{\tau(x)\tau(y)\to\tau(z)}(\ip{x}{y})\\
 &-2\big(K_x(y,z)+K_y(x,z)+K_z(x,y)\big)-c_T .
\end{aligned}
\end{equation}
It depends only on the pattern of $T$ and on the three inner products of $T$.

\begin{proposition}[fundamental identity]\label{prop:identity}
For every configuration $X$ of $14$ points carrying types with the prescribed cardinalities,
\[
  E(X)-e=\sum_{T}\lambda_T+\sum_{x\in X}\ \sum_{y,z\in X}K_x(y,z)\ \ge\ \sum_T\lambda_T ,
\]
the first sum running over the $\binom{14}{3}=364$ triples.
\end{proposition}
\begin{proof}
Split $\sum_x\sum_{y,z}K_x(y,z)$ according to the coincidences among $x,y,z$. Terms with $x,y,z$ distinct contribute
$2\sum_T\big(K_x(y,z)+K_y(x,z)+K_z(x,y)\big)$ ($K_x(y,z)=K_x(z,y)$ by symmetry of $F$). Terms with exactly two coincident
points: $y=z\neq x$ gives $\sum_k T(u)^\top F^{(k)}_{\tau(y)\tau(y)}T(u)(1-u^2)^k$ since $Q_k(u,u,1)=(1-u^2)^k$;
$y=x\neq z$ or $z=x\neq y$ give $\sum_{i,j}F^{(0)}_{(X,i),(\tau(z),j)}T_i(1)T_j(v)$ each, since $Q_k(1,v,v)=0$ for $k\ge1$.
Together these give $\sum_{\{x,y\}}\big(g_{\tau(x)}(\ip xy;\tau(y))+g_{\tau(y)}(\ip xy;\tau(x))\big)$. Terms $x=y=z$ give
$\sum_x\mathrm{sing}_{\tau(x)}$. Hence
\[
 \sum_{\{x,y\}}\varphi(\ip xy)=\sum_{\{x,y\}}\Psi(\ip xy)+\sum_x\sum_{y,z}K_x(y,z)-2\sum_T\sum_{\text{roots}}K-\sum_x\mathrm{sing}_{\tau(x)} .
\]
Every pair $\{x,y\}$ of types $\{a,b\}$ lies in exactly $m_c(ab)$ triples whose third vertex has type $c$; by
\eqref{eq:sharesum}, $\Psi_{ab}(\ip xy)=\sum_{T\ni x,y}\psi_{ab\to\tau(T\setminus\{x,y\})}(\ip xy)$. Summing over pairs and
regrouping by triples, and using \eqref{eq:constsum}, gives the identity. The inequality is Lemma~\ref{lem:pd}.
\end{proof}

\begin{corollary}\label{cor:bound}
Let $\mathcal D_{\pi}\subset[-1,1]^3$ be sets containing all inner-product types $(u,v,t)$ that triples of pattern $\pi$
can have in the class of configurations under consideration. If each slack function is nonnegative on the corresponding
$\mathcal D_\pi$, then $E(X)\ge e$ for all configurations in the class.
\end{corollary}

\paragraph{The typings used.}
\begin{itemize}
\item \emph{Untyped} ($T=1$): used in Case I; here we use the symmetric variant of~\cite{CohnWoo} (see Section~\ref{sec:case1}).
\item \emph{P/R typing}: $P$ = the minimal pair ($n_P=2$), $R$ = the other $12$ points. Patterns $PPR$ ($12$ triples),
$PRR$ ($132$), $RRR$ ($220$). With the shares $\psi_{PP\to R}=\Psi_{PP}/12$,
$\psi_{PR\to P}=\psi_{Ba}$, $\psi_{PR\to R}=(\Psi_{PR}-\psi_{Ba})/11$, $\psi_{RR\to P}=\psi_{Cb}$,
$\psi_{RR\to R}=(\Psi_{RR}-2\psi_{Cb})/10$, equation~\eqref{eq:sharesum} holds for each pair class.
If $m\in[a,b]$ then: $PPR$ types $(u,v,t)=(\ip PQ,\ip Pr,\ip Qr)$ satisfy $u\in[a,b]$, $v,t\ge u$; all other inner
products are $\ge a$.
\end{itemize}

%======================================================================
\section{Rigorous verification of slack positivity}\label{sec:verif}
%======================================================================
The four files named \texttt{certificate.json} in the directories listed in Appendix~\ref{sec:repro}
are the exact certificate data. They specify the energy offset, kernel matrices, polynomial coefficients and
Coulomb weights as rational numbers. The cap file also contains its factorizations and rational node centres.
These fixed data are verified as follows:
\begin{enumerate}[label=(\alph*)]
\item In Cases 0--II, every kernel matrix is proved positive definite by an exact $LDL^\top$
test over $\Q$. For Case III, the exact factorizations in Section~\ref{sec:capcert} prove positive semidefiniteness.
\item The program \path{audit_identity.py} checks the energy identity coefficient by coefficient over $\Q$,
including constants, pair multiplicities and the weights of $\varphi$. No floating-point reconstruction
of the kernel matrices is needed.
\item The slack functions are written exactly as $\lambda=\sum_q w_q\varphi(x_q)+P(u,v,t)$ with rational weights $w_q$ and a
polynomial $P\in\Q[u,v,t]$ computed in exact rational arithmetic (by expanding the definitions of
Section~\ref{sec:bounds}); exact invariance of $P$ under the relabelling symmetries of the pattern is checked, which allows
us to restrict to a fundamental domain.
\item Nonnegativity of $\lambda$ on its domain (or $\lambda\ge\theta$ for a threshold $\theta>0$ outside prescribed balls,
see Section~\ref{sec:cap}) is proved by the branch-and-bound program \path{bnb_global4} described next.
\end{enumerate}
The untyped certificates in Cases 0 and I use the normalization $4(E-e)$; the typed certificates use $E-e$.
The identity checker verifies these normalizations independently. Numerical optimization was used only to propose
the data, and is not part of the verification.

\paragraph{The branch-and-bound verifier.}
The domain is a union of root cells $\prod_i[c_i-r_i,c_i+r_i]$ (cells which provably miss the Gram set
$\{G\ge0\}$, $G=1+2uvt-u^2-v^2-t^2$, or violate a linear constraint are dropped by an exact rational test).
For every root cell the Taylor coefficients of $P$ at the cell centre $c$ (a double precision number, hence an exact
rational) are computed exactly over $\Q$ and rounded outward to intervals.  All further computations use interval
arithmetic with outward rounding after every operation (the program is compiled without value-changing floating-point
optimisations).  A cell is processed as follows.
\begin{enumerate}[label=(\roman*)]
\item The polynomial part is bounded on $|z_i|\le r_i$ ($z=x-c$) by $p_0+\sum_{\alpha\ne0}|p_\alpha|\,r^\alpha[-1,1]$.
\item The function $\varphi$ is expanded as $\varphi(c_q+z_q)=\varphi(c_q)\sum_{k\le D}\beta_k\big(z_q/(1-c_q)\big)^k$,
$\beta_k=\binom{2k}{k}4^{-k}\le1$, with the tail bounded by $\varphi(c_q)\varrho^{D+1}/(1-\varrho)$, $\varrho=r_q/(1-c_q)$,
if $\varrho<\tfrac12$; otherwise, using $w_q>0$ and the monotonicity of $\varphi$, by the constant $w_q\varphi(c_q-r_q)$.
The coefficients are enclosed by the interval recurrence $\beta_0=1$,
$\beta_k=\beta_{k-1}(2k-1)/(2k)$; no heuristic allowance for accumulated rounding is used.
\item If $G$ or a linear constraint $\ell\ge0$ of the domain changes sign on the cell, the lower bound is computed for
$\lambda-\sum_k\mu_kg_k$ ($g_k\in\{G,\ell\}$, $\mu_k\ge0$), which is $\le\lambda$ on the domain (S-procedure). The
multipliers are chosen heuristically (least squares on the gradients at the centre, followed by a coordinate ascent of the
resulting lower bound over the break points of this concave piecewise linear function); any choice $\mu\ge0$ is valid.
\item A cell is accepted if the lower bound is $\ge\theta$, discarded if $G<0$ or some $\ell<0$ on the whole cell, or if it
lies inside one of the excluded balls; otherwise it is bisected along its longest side, and the Taylor data of the
children are obtained from those of the parent by an interval Taylor shift by the exact double $\pm h$.
\item The \emph{true} centre of a child is the real number $c\pm h$, whose double precision representation may be rounded.
The rounding error is computed exactly (TwoSum) and accumulated, rounded upwards, in an error radius $e$; all geometric tests
use the half widths $r+e$, and all quantities evaluated at the centre (Gram polynomial, $\varphi$-coefficients, values of
the linear constraints) are evaluated on the interval $[c-e,c+e]$. (In all our runs $e\le 4\cdot10^{-16}$.)
\end{enumerate}
The run succeeds if every cell is accepted or discarded before its width drops below a prescribed minimum.
A second program, \path{bnb_local3}, works in the same way in node-adapted coordinates around the nodes of $\Xs$; it is
described in Section~\ref{sec:cap}.

\paragraph{Implementation and audit.}
All current global results use \path{bnb_global4}; the local calculations use \path{bnb_local3}.
Both traverse the tree depth first using OpenMP tasks, with a fixed depth limit. A depth or resource limit, an unresolved
cell, an input parsing failure, or an early abort prevents a verified result. The programs require IEEE binary64 arithmetic in
round-to-nearest mode and are compiled with \texttt{-ffp-contract=off -fno-fast-math}.
Each elementary arithmetic interval is enlarged using \texttt{nextafter}; square roots and reciprocals are likewise
enclosed. Degrees are at most $40$, so the integer binomial coefficients used in Taylor shifts are exactly representable.
Every discarded node ball is covered by an independently checked local region.
The complete source of both verifiers is included in \path{Certificates/common/}.

%======================================================================
\section{Case I: a lower bound on every inner product}\label{sec:case1}
%======================================================================
For the untyped bound we use symmetric three-point kernels, as in~\cite{CohnWoo}: for
$k=0,\dots,K$ let $F_k\PSD$ be indexed by $0\le i,j\le K-k$, put $S_k(u,v,t)=\big(T_i(u)T_j(v)Q_k(u,v,t)\big)_{ij}$ and
\[
 R_k(u,v,t)=\sum_{\pi\in S_3}S_k(\pi(u,v,t))+\frac{1}{N-2}\sum_{w\in\{u,v,t\}}\,\sum_{\pi\in S_3}S_k(\pi(w,w,1))
 +\frac{6\,\delta_{k0}}{(N-1)(N-2)}J .
\]
The slack of a triple with inner products $(u,v,t)$ is
\[
 \lambda(u,v,t)=\tfrac13\big(\varphi(u)+\varphi(v)+\varphi(t)\big)-\frac{e}{\binom N2}-\sum_k\ip{F_k}{R_k(u,v,t)},
\]
and exactly as in Proposition~\ref{prop:identity} one obtains
$4(E(X)-e)=\sum_T\lambda_T+(\text{PSD terms})\ge\sum_T\lambda_T$. Hence:

\begin{proposition}\label{prop:caseI}
If $\lambda\ge0$ on $\{(u,v,t)\in[a,1)^3:\ 1+2uvt-u^2-v^2-t^2\ge0\}$, then $E(X)\ge e$ for every configuration $X$ all of whose
inner products are $\ge a$.
\end{proposition}

\paragraph{Certificate.} The degree-$40$ certificate in \path{CaseI_a-0.98_D40/certificate.json}
has the following verified bounds; the decimal energy and gap in the table are approximate:
\begin{center}
\begin{tabular}{llllr}
\toprule
$a$ & degree & certified $e$ & $e-\Es$ & cells (failures)\\
\midrule
$-0.98$ & $40$ & $69.30715094\ldots$ & $+7.877\cdot10^{-4}$ & $376\,059$ ($0$)\\
\bottomrule
\end{tabular}
\end{center}
The exact offset is $e_{\rm I}=163646883327020528602849/2361183241434822606848$.
The positivity verification uses \path{bnb_global4}.

\subsection{A priori separation of the points}\label{sec:sep}
The untyped slack also bounds the largest inner product of any competitor.
The cutoff $t_0=0.8$ uses the numerical convention in Section~1.

\begin{proposition}[separation]\label{prop:sep}
Every configuration $X$ of $14$ distinct points on $\Sph$ which contains two points with $\ip{x_i}{x_j}\ge t_0=0.8$
satisfies $E(X)\ge\Es+0.086$. Consequently all inner products of a configuration with $E(X)\le\Es$ are smaller than $t_0$.
\end{proposition}
\begin{proof}
Let $(e,F_0,\dots,F_K)$ be an untyped certificate whose slack satisfies
\begin{equation}\label{eq:sep}
 \lambda\ge0\ \text{ on }\{(u,v,t)\in[-1,1)^3:\ G\ge0\},\qquad
 \lambda\ge c\ \text{ on the subset where }\max(u,v,t)\ge t_0 .
\end{equation}
If $\ip{x_i}{x_j}\ge t_0$, then each of the $12$ triples containing $x_i$ and $x_j$ has slack $\ge c$, and all other
slacks are nonnegative. Since $4(E(X)-e)\ge\sum_T\lambda_T$ (the identity preceding Proposition~\ref{prop:caseI}), we get
$E(X)\ge e+3c$. The degree-$16$ certificate in
\path{Case0_separation_T0.8_D16/certificate.json} has
\[
 e_0=\tfrac{5110096645419161624321}{73786976294838206464}=69.25472355\ldots=\Es-0.05163\ldots,\qquad c=0.0459 .
\]
By the exact $S_3$-symmetry of the slack it suffices to consider $u\le v\le t$, where the second set is $\{t\ge t_0\}$.
\path{bnb_global4} verifies both parts of \eqref{eq:sep}: $29\,287$ cells for $\lambda\ge0$ and $10\,041$ cells for
$\lambda\ge c$ on $\{t\ge t_0\}$, without failures. Hence $E(X)\ge e_0+3c=69.39242\ldots>\Es+0.086$.
\end{proof}
Proposition~\ref{prop:sep} allows us to intersect all slack
domains of Case~III with $\{u,v,t\le t_0\}$; this removes the region of nearly coincident points, where the kernel
polynomials have to follow the singularity of $\varphi$.

%======================================================================
\section{Case II: the intermediate slab}\label{sec:slabs}
%======================================================================
We use the $P/R$ typing with shares of Section~\ref{sec:bounds} with $K=10$ and polynomial shares of degree $20$.
For $m\in[-0.99,-0.98]$ the slack domains are (Corollary~\ref{cor:bound}):
$PPR$: $u\in[-0.99,-0.98]$, $v,t\ge u$; $PRR$ and $RRR$: all inner products $\ge-0.99$.
By the exact invariance of the slack polynomials (checked) we may restrict to $t\ge v$ for $PPR$, $v\ge u$ for $PRR$ and
$u\le v\le t$ for $RRR$. The exact certificate in \path{Slab_m099_m098/certificate.json} has
\[
 e_{\rm II}=\tfrac{20455735517135044064649}{295147905179352825856}=69.30672777\ldots,\qquad e_{\rm II}-\Es\approx3.645\cdot10^{-4},
\]
and \path{bnb_global4} verifies nonnegativity of the three slacks ($53\,303$, $219\,381$ and $76\,285$ cells for $PPR$, $PRR$, $RRR$; no failures).
Together with Case~I (certificate for $a=-0.98$) this excludes all configurations with $m\ge-0.99$.

%======================================================================
\section{Case III: the nearly antipodal cap}\label{sec:cap}
%======================================================================
Throughout this section $X$ is a configuration with $m(X)=\ip PQ\in[-1,-0.99]$, $\{P,Q\}$ a minimal pair, all of whose
inner products are smaller than $t_0=0.8$ (Proposition~\ref{prop:sep}), and we use the $P/R$ typing of
Section~\ref{sec:bounds} with $K=10$ (degree $20$). For every point $x\ne P,Q$ we have
$\ip Px=\ip{P+Q}{x}-\ip Qx>-|P+Q|-t_0>-0.1415-t_0>-c_1$ with $c_1:=0.9416$ (note $|P+Q|=\sqrt{2+2m}<0.1415$), and
likewise for $Q$. The slack domains
are therefore
\begin{align*}
 \mathcal D_{PPR}&=\{u\in[-1,-0.99],\ v,t\in[-c_1,t_0],\ v,t\ge u,\ G\ge0\},\\
 \mathcal D_{PRR}&=\{u,v\in[-c_1,t_0],\ t\in[-1,t_0],\ G\ge0\},\qquad
 \mathcal D_{RRR}=\{(u,v,t)\in[-1,t_0]^3:\ G\ge0\}.
\end{align*}
Since $\Xs$ belongs to this class, a nonnegative-slack certificate cannot have $e>\Es$.
We use $e<\Es$ together with quantitative slack bounds to localize a competitor, then apply local minimality.

\subsection{Nodes and class values}
The inner-product types $s^\ast=(u,v,t)$ of the triples of $\Xs$, in all orderings compatible with the pattern, are the
\emph{nodes} of the pattern; $\mathcal N_\pi$ denotes the set of nodes of the pattern $\pi$ and
\[
 d_\pi(s)=\min_{s^\ast\in\mathcal N_\pi}|s-s^\ast|_2\qquad(s\in\R^3).
\]
With $h=\hs$ and $G_=$, $G_{\ne}$ as in Section~\ref{sec:rigidity}: $\mathcal N_{PPR}=\{(-1,h,-h),(-1,-h,h)\}$;
$\mathcal N_{PRR}$ consists of the $12$ points $(\sigma h,\sigma'h,g)$ with $\sigma,\sigma'\in\{\pm1\}$, $g\in G_=$ if
$\sigma=\sigma'$ and $g\in G_{\ne}$ otherwise; $\mathcal N_{RRR}$ consists of the $55$ ordered types of three ring points of
$\Xs$. The $PPR$ nodes lie on the edge $\{u=-1\}\cap\{G=0\}$ of the domain, the nodes $(\pm h,\pm h,h^2-r^2)$ of $PRR$ lie
on the Gram surface $\{G=0\}$ (the three points are coplanar with the centre); all other nodes are interior points.

The pairs of points of $X$ fall into three \emph{pair classes} with the \emph{class values}
\[
 V_{PP}=\{-1\},\qquad V_{PR}=\{h,-h\},\qquad V_{RR}=G_=\cup G_{\ne}
\]
(the inner products of the corresponding pairs of $\Xs$). Every coordinate of a node is a class value of the pair class
of that coordinate. Consequently
\begin{equation}\label{eq:coord}
 \begin{gathered}
 \text{if a coordinate of $s$ is at distance $\ge\delta$ from its pair-class values,}\\
 \text{then }d_\pi(s)\ge\delta.
 \end{gathered}
\end{equation}

\subsection{Three analytic reductions}
\begin{lemma}[counting]\label{lem:count}
Let $(e,\lambda)$ be a $P/R$-typed certificate with $\lambda_\pi\ge0$ on $\mathcal D_\pi$ for all patterns, and let
$\tau,\theta>0$ be such that $\lambda_\pi(s)\ge\theta$ for all $s\in\mathcal D_\pi$ with $d_\pi(s)\ge\tau$. If
$E(X)-e<12\,\theta$, then every inner product of $X$ lies within $\tau$ of a class value of its pair class.
\end{lemma}
\begin{proof}
Suppose that the inner product of a pair $\{x,y\}$ has distance $\ge\tau$ from all class values of its class. The pair lies
in exactly $12$ triples $\{x,y,z\}$, and by \eqref{eq:coord} each of them has a type with $d_\pi\ge\tau$, hence a slack
$\ge\theta$. All other slacks are nonnegative, so Proposition~\ref{prop:identity} gives
$E(X)-e\ge\sum_T\lambda_T\ge12\,\theta$.
\end{proof}
Thus the threshold which defines the tubes around the nodes only has to exceed $(E(X)-e)/12$ instead of $E(X)-e$: the
energy excess is shared by the twelve triples through a displaced pair. Moreover the conclusion concerns pairs, which
is what Lemma~\ref{lem:rigidity} requires. (The same device was used in Proposition~\ref{prop:sep}.)

\begin{lemma}[sign coupling]\label{lem:sign}
Let $0<\tau\le10^{-3}$, $|\ip PQ+1|\le\tau$, and let $x\in\Sph$ satisfy $|\ip Px-\sigma h|\le\tau$ and
$|\ip Qx-\sigma'h|\le\tau$ with $\sigma,\sigma'\in\{\pm1\}$. Then $\sigma'=-\sigma$.
\end{lemma}
\begin{proof}
$|\ip Px+\ip Qx|=|\ip{P+Q}{x}|\le|P+Q|=\sqrt{2+2\ip PQ}\le\sqrt{2\tau}<0.05$, whereas $\sigma'=\sigma$ would give
$|\ip Px+\ip Qx|\ge2h-2\tau>0.9$.
\end{proof}

\begin{lemma}[class coupling]\label{lem:class}
Let $\gamma=\min\{|g-g'|:\ g\in G_=,\ g'\in G_{\ne}\}>0.018$, as verified in ball arithmetic,
and let $0<\tau\le10^{-3}$. Let $x,y$
satisfy $|\ip Px-\sigma h|\le\tau$, $|\ip Py-\sigma'h|\le\tau$ and $|\ip xy-g|\le\tau$, where $g\in G_=\cup G_{\ne}$ is of
the \emph{wrong} class ($g\in G_{\ne}$ if $\sigma=\sigma'$, $g\in G_=$ otherwise). Then the type $s$ of the triple
$\{P,x,y\}$ satisfies $d_{PRR}(s)\ge\gamma-\tau>0.017$.
\end{lemma}
\begin{proof}
A node $(\varrho h,\varrho'h,g^\ast)$ with $(\varrho,\varrho')\ne(\sigma,\sigma')$ has distance $\ge2h-\tau$ from $s$. A node
with $(\varrho,\varrho')=(\sigma,\sigma')$ has $g^\ast$ in the right class, hence
$|\ip xy-g^\ast|\ge|g-g^\ast|-\tau\ge\gamma-\tau$.
\end{proof}

\subsection{Sufficient certificate inequalities}
\begin{proposition}\label{prop:cap}
Let $(e,\lambda)$ be a $P/R$-typed certificate with $e\le\Es$ and let $\tau,\theta_{\rm loc},\theta_{\rm far},\rho_1>0$ be
such that for every pattern $\pi$
\begin{align}
 \lambda_\pi&\ge0 &&\text{on }\mathcal D_\pi, \label{eq:c1}\\
 \lambda_\pi&\ge\theta_{\rm loc} &&\text{on }\{s\in\mathcal D_\pi:\ d_\pi(s)\ge\tau\}, \label{eq:c2}\\
 \lambda_\pi&\ge\theta_{\rm far} &&\text{on }\{s\in\mathcal D_\pi:\ d_\pi(s)\ge\rho_1\}, \label{eq:c3}
\end{align}
where $12\,\theta_{\rm loc}>\Es-e$, $\theta_{\rm far}>\Es-e$, $\tau\le10^{-3}$ and $\rho_1\le0.017$. Then every configuration
$X$ of the class of this section with $E(X)\le\Es$ satisfies the hypotheses (i)--(iii) of Lemma~\ref{lem:rigidity} with
this $\tau$.
\end{proposition}
\begin{proof}
By Lemma~\ref{lem:count} (with $\theta=\theta_{\rm loc}$) every inner product of $X$ is within $\tau$ of a class value of
its class. Hence $|\ip PQ+1|\le\tau$, which is (i), and for every $x_i$ there are signs with
$|\ip P{x_i}-\sigma_ih|\le\tau$, $|\ip Q{x_i}-\sigma_i'h|\le\tau$; by Lemma~\ref{lem:sign}, $\sigma_i'=-\sigma_i$, which is
(ii). For $i\ne j$ the inner product $\ip{x_i}{x_j}$ is within $\tau$ of some $g\in G_=\cup G_{\ne}$. If $g$ were of the wrong
class, the triple $T=\{P,x_i,x_j\}$ would satisfy $d_{PRR}\ge0.017\ge\rho_1$ by Lemma~\ref{lem:class}, hence by
\eqref{eq:c3}, \eqref{eq:c1} and Proposition~\ref{prop:identity}
\[
 \theta_{\rm far}\le\lambda_T\le\sum_{T'}\lambda_{T'}\le E(X)-e\le\Es-e<\theta_{\rm far},
\]
a contradiction. This gives (iii).
\end{proof}
Only inner products, i.e.\ the patterns $PPR$ and $PRR$, enter the rigidity lemma; the $RRR$ slack is needed for
\eqref{eq:c1} and for the counting.



\subsection{The certificate}\label{sec:capcert}
\paragraph{Construction and exact data.}
The cap data are in \path{CaseIII_cap_D20/certificate.json}. The search used facial reduction,
$F=NGN^\top$, and a degree-$20$ conic optimization problem with sampled slack and
node-derivative constraints. These steps only produce candidate data. For verification, $N$ and $G$ are fixed rational
matrices, and the identities $F=NGN^\top$ and $G\succ0$ are checked exactly. No claim that the rational $N$ spans an
exact null space of the candidate moment matrix is needed.

Set $\delta_0=\DELTAZ$ and
\[
 e=\Es_{\rm lo}-364\delta_0,\qquad
 \Es_{\rm lo}=69306363296626421258489991276382007864282129573758\cdot10^{-48}.
\]
The enclosure of Lemma~\ref{lem:hstar} proves $\Es_{\rm lo}\le\Es$ and
\[
 e=\frac{34653181648313028629244995638191003932141064786879}{5\cdot10^{47}},
 \qquad 0<\Es-e<3.65\cdot10^{-13}.
\]
The stored slacks have the exact form $\sum_q w_q\varphi(x_q)+P(u,v,t)$ with $w_q\in\Q_{>0}$ and
$P\in\Q[u,v,t]$. The program \path{audit_identity.py} checks the fundamental energy identity
coefficient by coefficient over $\Q$, directly from the stored kernel matrices and slack polynomials.
It also checks every relabelling symmetry, including the singular weights. Thus one may restrict to
$t\ge v$ for $PPR$, $v\ge u$ for $PRR$, and $u\le v\le t$ for $RRR$.

For each node we use a rational approximation $\tilde s$. The program \path{audit_geometry.py} enumerates the triple
types of the two regular hexagons and the poles using their exact symbolic pair classes, and verifies in ball arithmetic
that the archived list contains all $2+12+55$ nodes and that $|\tilde s-s^\ast|<10^{-50}$.
The stored centres and matrices are rational approximations; no exact interpolation or stationarity
condition at these centres is assumed. The verification below uses the actual rational
Taylor coefficients and explicitly bounds their nonzero constant and linear terms.

\paragraph{Local coordinates and coverage.}
For each fundamental-domain node, the cap certificate supplies a rational matrix $M$ and positive rational
numbers $s_1,s_2,s_3$, together with the norm bound and multipliers used below. These are stored in the
\path{local_coordinates} field. Numerical diagonalization was used to propose these data;
reproduction uses the stored rational values.
Define
\[
 x=\tilde s+My,\qquad
 \mathcal E=\left\{y:\sum_i y_i^2/s_i^2\le1\right\},\qquad
 R_{\rm out}=0.01,\quad \rho_g=0.0095.
\]
All matrix inequalities subsequently used are checked over $\Q$. With $S=\operatorname{diag}(s_1,s_2,s_3)$,
positive definiteness of
\[
 R_{\rm out}^2 S^{-2}-M^\top M,\qquad
 M^\top M-(\rho_g+10^{-14})^2S^{-2},\qquad
 c_M^2 I-M^\top M
\]
proves, respectively, that $|My|\le R_{\rm out}$ on $\mathcal E$, that
$M\mathcal E$ contains $B(0,\rho_g+10^{-14})$, and that $\|M\|_2\le c_M$.
The global centres differ from $\tilde s$ by less than $10^{-16}$, and the stored radii differ from $\rho_g$ by less
than $10^{-16}$, as checked exactly. Thus each omitted ball is contained in
$B(\tilde s,\rho_g+10^{-14})$ and is covered by the local ellipsoid.

Expanding each $\varphi$ term to degree $20$ at $\tilde s$, enclosing its coefficients in ball arithmetic, and bounding
the geometric series tail produces an exact rational polynomial $\Lambda(y)$ and rational $\varepsilon>0$ such that
\[
 |\lambda(\tilde s+My)-\Lambda(y)|\le\varepsilon
 \quad\hbox{whenever }|My|\le R_{\rm out}.
\]
Polynomial substitutions and Taylor shifts are exact rational operations. The roots of the local branch-and-bound
calculation cover the full bounding box of $\mathcal E$. Affine domain cuts are rounded conservatively, so the tested
domain contains the true domain.

\paragraph{The inner node estimate.}
Choose rational multipliers $m_G,\nu\ge0$, with $\nu=0$ except at the edge node, and put
\[
 L(y)=\Lambda(y)-m_G G(\tilde s+My)-\nu(u+1)
     =c_0+g\cdot y+\tfrac12 y^\top H_Ly+R(y).
\]
The subtracted functions are nonnegative on the true domain. Exact positive-definiteness tests give
$H_L\succeq(9/10)I$. If $R(y)=\sum_{|\alpha|\ge3}r_\alpha y^\alpha$, the bound
\[
 |R(y)|\le M_r|y|^3,\qquad
 M_r=\sum_{|\alpha|\ge3}|r_\alpha|y_{\rm in}^{|\alpha|-3},
 \qquad M_r y_{\rm in}\le\frac9{40}
\]
is certified in ball arithmetic; $c_M y_{\rm in}\le R_{\rm out}$ is checked exactly.
Thus, for $|y|\le y_{\rm in}$ in the domain,
\[
 \lambda\ge c_0-\varepsilon-|g|\,|y|+\frac9{40}|y|^2.
\]
The exact rational test
$c_0-\varepsilon\ge |g|^2/(9/20)$ is stronger than the condition obtained by minimizing this quadratic and proves
$\lambda\ge0$. Moreover $\lambda\ge\theta_{\rm loc}$ whenever
\[
 y_{\rm tin}:=\frac{|g|}{9/40}
       +\sqrt{\frac{\theta_{\rm loc}+\varepsilon}{9/40}}
       \le |y|\le y_{\rm in}.
\]
The computation uses a rational upper bound for $y_{\rm tin}$; for every node it is at most
$4.22\cdot10^{-7}<y_{\rm in}$.

\paragraph{The remaining local and global regions.}
The local interval verifier proves nonnegativity on the remainder of $\mathcal E$ intersected with the domain.
It proves the stronger bound $\lambda\ge\theta_{\rm loc}$ except in cells whose norm is at most the reported
$y_{\rm tube}$. This number is an outward bound, exported with enough digits to recover the binary number exactly;
its conversion and use in the final inequality are checked over $\Q$.
Combining this calculation with the inner node estimate proves nonnegativity everywhere in the local region and
the threshold bound outside distance
\[
 c_M\max(y_{\rm tube},y_{\rm tin})+10^{-50}
\]
from the corresponding true node.

The global verifier proves $\lambda\ge\theta_{\rm far}$ outside the omitted balls. These balls are covered by the local
calculations, as just shown. Every omitted ball lies within distance $0.01$ of its corresponding true node.
We therefore obtain~\eqref{eq:c1}--\eqref{eq:c3} with $\rho_1=0.01$, provided that the preceding local tube radii are
at most $\tau$. The choice $0.01$ leaves a strict margin for both rational-to-double centre conversion and the
approximation of the true nodes.

\subsection{Result}
With $\theta_{\rm loc}=\THLOC>(\Es-e)/12$ and $\theta_{\rm far}=\THFAR>\Es-e$:
\begin{center}
\begin{tabular}{lrrr}
\toprule
 & nodes / root cells & cells & failures\\
\midrule
local runs ($22$ nodes of the fundamental domains) & $22\times64$ & $498\,434$ & $0$\\
global run $PPR$ & $20$ & \NPPR & $0$\\
global run $PRR$ & $427$ & \NPRR & $0$\\
global run $RRR$ & $229$ & \NRRR & $0$\\
\bottomrule
\end{tabular}
\end{center}
There are $22$ nodes in the fundamental domains: $1$ of $PPR$, $9$ of $PRR$, and $12$ of $RRR$.
The node estimate holds at all of them. The exact and interval checks give
$c_M\le\CMAX$, $\varepsilon\le10^{-34}$, $y_{\rm tube}<8.211\cdot10^{-7}$, and
\[
 \tau=\max_{\rm nodes}c_M\max(y_{\rm tube},y_{\rm tin})+10^{-50}<\TUBE .
\]
The global runs use target root-cell width approximately $0.2$ and minimum half-width approximately $10^{-10}$;
the smallest accepted lower bounds are approximately
$2.3\cdot10^{-12}$, $1.2\cdot10^{-12}$ and $1.0077\cdot10^{-12}$ (threshold $10^{-12}$).
\begin{proof}[Proof of Theorem~\ref{thm:C}]
The certificate satisfies \eqref{eq:c1}--\eqref{eq:c3} with $\tau=\TUBE\le10^{-3}$, $\rho_1=0.01<0.017$,
$12\,\theta_{\rm loc}>4.79\cdot10^{-13}>3.65\cdot10^{-13}>\Es-e$ and $\theta_{\rm far}>9.9\cdot10^{-13}>\Es-e$;
Proposition~\ref{prop:cap} applies.
\end{proof}

%======================================================================
\section{Rigidity: from inner-product types to coordinates}\label{sec:rigidity}
%======================================================================
Write $h=\hs$, $r=\sqrt{1-h^2}$ and
\[
 G_{=}=\{h^2+\tfrac{r^2}{2},\ h^2-\tfrac{r^2}{2},\ h^2-r^2\},\qquad
 G_{\neq}=\{-h^2+\tfrac{\sqrt3}{2}r^2,\ -h^2,\ -h^2-\tfrac{\sqrt3}{2}r^2\}
\]
(the inner products of two ring points of $\Xs$ in the same ring, resp.\ in different rings).

\begin{lemma}[rigidity]\label{lem:rigidity}
Let $0<\tau\le10^{-3}$ and let $X=\{P,Q,x_1,\dots,x_{12}\}\subset\Sph$ satisfy
\begin{enumerate}[label=(\roman*)]
\item $|\ip PQ+1|\le\tau$;
\item for every $i$ there is $\sigma_i\in\{\pm1\}$ with $|\ip P{x_i}-\sigma_ih|\le\tau$ and $|\ip Q{x_i}+\sigma_ih|\le\tau$;
\item for all $i\neq j$ there is $g\in G_{=}$ (if $\sigma_i=\sigma_j$), resp.\ $g\in G_{\neq}$ (if $\sigma_i\neq\sigma_j$), with
      $|\ip{x_i}{x_j}-g|\le\tau$.
\end{enumerate}
Then $\#\{i:\sigma_i=1\}=\#\{i:\sigma_i=-1\}=6$, and after an orthogonal change of coordinates and a relabelling of the points
the gauge coordinates $p$ of $X$ (Section~\ref{sec:local}) satisfy $|p-p^\ast|_\infty\le 9\tau$ and $|p-p^\ast|_2\le 23\tau$.
\end{lemma}
\begin{proof}
Put $P=e_3$ by an orthogonal map and write
$x_i=(\rho_i\cos\theta_i,\rho_i\sin\theta_i,z_i)$, where $\rho_i=\sqrt{1-z_i^2}$.
Hypothesis (ii) gives $|z_i-\sigma_i h|\le\tau$ and
\[
 |\rho_i-r|\le\beta\tau,\qquad
 \beta=\frac{h+\tau}{\sqrt{1-(h+\tau)^2}}<0.513.
\]
For points with equal signs and $g=h^2+r^2c$ from (iii),
\[
 \cos(\theta_i-\theta_j)-c
 =\frac{(\langle x_i,x_j\rangle-g)-(z_i z_j-h^2)-c(\rho_i\rho_j-r^2)}{\rho_i\rho_j}.
\]
Here $|z_i z_j-h^2|\le2h\tau+\tau^2$,
$|\rho_i\rho_j-r^2|\le2r\beta\tau+\beta^2\tau^2$, and
$\rho_i\rho_j\ge(r-\beta\tau)^2$. Therefore
\[
 |\cos(\theta_i-\theta_j)-c|\le3.569\tau
 \quad(c\in\{1/2,-1/2,-1\}).
\]
For opposite signs the same calculation, with $c\in\{\sqrt3/2,0,-\sqrt3/2\}$ and
$z_i z_j$ compared to $-h^2$, gives the bound $3.414\tau$.

Set $a=1.16(3.569)$, $b=1.001(3.414)$, and $c_0=2.32(3.414)$, all exact decimals.
The mean value theorem for cosine shows that an equal-sign angular distance is within $a\tau$ of
$\pi/3$ or $2\pi/3$, or belongs to the separate option near $\pi$.
An opposite-sign distance is within $b\tau$ of $\pi/2$ or within $c_0\tau$ of $\pi/6$ or $5\pi/6$.
For example, the first bound follows from
$1.16\sin(\pi/3-a\tau)>1$; the other two follow from
$1.001\cos(b\tau)>1$ and $2.32\sin(\pi/6-c_0\tau)>1$.
These inequalities and the separation of the three options are checked for $\tau\le10^{-3}$ by
\path{audit_geometry.py}, using rational constants and ball arithmetic.

Every equal-sign angular distance is at least $\pi/3-a\tau$.
Since $7(\pi/3-a\tau)>2\pi$, each sign occurs at most six times and hence exactly six times.
In circular order, all six consecutive gaps are at least $\pi/3-a\tau$ and sum to $2\pi$.
Each is consequently at most $\pi/3+5a\tau$.
This upper bound excludes both the $2\pi/3$ option and the near-$\pi$ option, so every gap is within $a\tau$ of $\pi/3$.

Choose an upper-ring reference point with azimuth zero and label the remaining upper points in circular order.
The gap bounds first put them within $\min(k,6-k)a\tau$ of $k\pi/3$.
Direct comparison with the reference point then gives error at most $a\tau$ for $k=1,2,4,5$; the angular options are
too far apart to be confused. Comparing the opposite point with the point at $k=1$ bounds its error by $2a\tau$.

Comparison of a lower-ring point with the reference point puts its angle within $c_0\tau$ of an odd multiple
of $\pi/6$. Of the two upper ideal vertices at angular difference $\pi/2$ from that location, at least one is not
the vertex opposite the reference point and therefore has error at most $a\tau$.
The coarse error bound $(c_0+a)\tau$ forces the comparison with that vertex to use the $\pi/2$ option.
The sharper bound for this option puts the lower point within $(a+b)\tau$ of its ideal location.
Two lower points cannot occupy the same ideal location, since
$2(a+b)\tau<\pi/3-a\tau$. There are six lower points and six such locations, so every location occurs once.
This proves the required $30^\circ$ rotation between the rings.

Write $Q=(A,B,-w)$ with $w=-\langle P,Q\rangle\in[1-\tau,1]$.
For $z_i=\sigma_i h+\epsilon_i$ and
$\langle Q,x_i\rangle=-\sigma_i h+\delta_i$, where $|\epsilon_i|,|\delta_i|\le\tau$,
\[
 |\rho_i(A\cos\theta_i+B\sin\theta_i)|
 =|\delta_i+w\epsilon_i+\sigma_i h(w-1)|\le(2+h)\tau.
\]
The upper reference point gives $|A|\le(2+h)\tau/(r-\beta\tau)$.
Using the lower point whose angle is within $(a+b)\tau$ of $\pi/2$ then bounds $|B|$.
Both bounds are less than $2.79\tau$, as certified by \path{audit_geometry.py}.
The twelve height errors, eleven free azimuth errors, and two pole-coordinate errors consequently satisfy
\[
 |p-p^\ast|_2^2\le
 \bigl(12+4a^2+(2a)^2+6(a+b)^2+2(2.79)^2\bigr)\tau^2
 <23^2\tau^2,
 \qquad |p-p^\ast|_\infty<9\tau.
\]
The rational and interval inequalities above hold uniformly on $0<\tau\le10^{-3}$: the error bounds increase with
$\tau$ and their denominators remain positive. Their evaluation at $10^{-3}$ therefore suffices.
\end{proof}

%======================================================================
\section{Local minimality}\label{sec:local}
%======================================================================
\paragraph{Gauge coordinates.} For configurations near $\Xs$ we fix $P=e_3$, write the second pole as $Q=(a,b,-\sqrt{1-a^2-b^2})$,
the first upper-ring point as $(\sqrt{1-z_2^2},0,z_2)$ (azimuth fixed to $0$) and all other ring points as
$(\sqrt{1-z_i^2}\cos\theta_i,\sqrt{1-z_i^2}\sin\theta_i,z_i)$; this gives $25$ parameters $p$, with $p^\ast$ the parameters of $\Xs$.
Every configuration whose points are within the bounds of Lemma~\ref{lem:rigidity} of $\Xs$ is, after an orthogonal map, of this form.

\begin{proposition}[quantitative local minimality]\label{prop:local}
For $0<|p-p^\ast|_2\le 0.0075$ one has $E(p)>E(p^\ast)=\Es$.
\end{proposition}
\begin{proof}
Let $\Delta=p-p^\ast$. Criticality from Lemma~\ref{lem:hstar} and Taylor's theorem give
\[
 E(p)-\Es=\tfrac12\Delta^\top H\Delta+C(\Delta)+R_4(\Delta).
\]
The program \path{audit_local.py} obtains interval Taylor coefficients by the truncated power-series arithmetic
of \path{localmin_jet.py}. It verifies the following bounds with rational matrix tests and ball arithmetic:
\[
 H\succeq0.0964I,\qquad |C(\Delta)|\le5.04|\Delta|_2^3,
 \qquad |R_4(\Delta)|\le147|\Delta|_2^4
 \quad\hbox{if }|\Delta|_\infty\le0.008.
\]
Here are the details of the norm estimates. For the Hessian, the exact rational midpoint matrix minus $0.0965I$
is positive definite, and the Frobenius norm of its error matrix is less than $0.0001$.
For the cubic form, its symmetric coefficient tensor is flattened to a $25\times625$ matrix $F$.
If $F_0$ is its rational midpoint, exact positive definiteness of
$(5.039)^2I-F_0F_0^\top$, together with an interval Frobenius bound less than $0.001$ for $F-F_0$,
gives the stated cubic bound.

Each of the $91$ pair energies depends on at most four gauge coordinates. Its fourth-order coefficient tensor is
enclosed over the entire box $|\Delta|_\infty\le0.008$, so the enclosure holds at every intermediate point in the
one-dimensional Taylor remainder along $p^\ast+t\Delta$.
For each pair, flattening gives a symmetric matrix with rational midpoint $M_{ij}$.
Exact positive-definiteness tests on $q_{ij}I\pm M_{ij}$, followed by an interval Frobenius bound for the error,
give $|R_{ij}|\le s_{ij}|\Delta_{ij}|_2^4$.
Here $\Delta_i$ denotes the parameter block belonging to point $i$, and
$|\Delta_{ij}|_2^2=|\Delta_i|_2^2+|\Delta_j|_2^2$.
Thus
\[
 |R_4|\le2\sum_i |\Delta_i|_2^4\sum_{j\ne i}s_{ij}
 \le\left(2\max_i\sum_{j\ne i}s_{ij}\right)|\Delta|_2^4
 \le147|\Delta|_2^4.
\]
The last inequality is checked in ball arithmetic. The rational proposals $q_{ij}$ are supplied in
\path{Local/local_norm_bounds.json}, relative to \path{Certificates/}.
They are accepted as bounds only after the exact matrix tests.

For $0<|\Delta|_2\le0.0075$, the box condition holds and
\[
 E(p)-\Es\ge |\Delta|_2^2
 \bigl(0.0482-5.04|\Delta|_2-147|\Delta|_2^2\bigr)>0.
\]
The bracket decreases with $|\Delta|_2$ and equals the positive rational number $341/160000$ at $0.0075$.
\end{proof}

%======================================================================
\appendix
\section{Certificates and reproduction}\label{sec:repro}
%======================================================================
Each case has an exact \texttt{certificate.json} file, together with problem files and compressed root data.
The JSON file specifies the rational matrices and slack coefficients.
The directories in the table are relative to \path{Certificates/}.
All verification code is in \path{Certificates/common/}.
\begin{center}
\begin{tabular}{lll}
\toprule
statement & certificate directory & verification runs\\
\midrule
Proposition~\ref{prop:sep} (Case 0) & \path{Case0_separation_T0.8_D16} & $2$ global\\
Proposition~\ref{prop:caseI} (Case I) & \path{CaseI_a-0.98_D40} & $1$ global\\
Case II, Section~\ref{sec:slabs} & \path{Slab_m099_m098} & $3$ global\\
Theorem~\ref{thm:C} (Case III) & \path{CaseIII_cap_D20} & $22$ local, $3$ global\\
\bottomrule
\end{tabular}
\end{center}
The problem/root file prefixes are \path{case0_pos}, \path{case0_close} (sharing \path{case0_roots}),
\texttt{v3t98}, \path{s1_D20_PPR}, \path{s1_D20_PRR}, and \path{s1_D20_RRR}.
Within the cap directory, the global prefixes are \path{global/g_PPR}, \path{global/g_PRR}, and
\path{global/g_RRR}. The $22$ local prefixes are \path{local/nXX}, with
\[
\begin{split}
\mathrm{XX}\in\{&00,02,03,04,05,06,07,11,12,13,14,15,16,17,19,\\
                 &25,26,34,36,44,48,56\}.
\end{split}
\]
Problem files have suffix \path{_ts.txt}; compressed root files have suffix \path{_roots.txt.xz}.
The complete data schema and the meanings of every problem-file field are documented in
\path{Certificates/FORMAT.md}. The top-level \texttt{SHA256SUMS} identifies all release files.
Hashes establish file identity; the mathematical verification additionally checks the identities,
arithmetic bounds and domain coverage described in this paper.

\paragraph{Analytic checks and their programs.}
Lemma~\ref{lem:hstar} uses \path{enclose_h.py}.
The exact kernel identities and matrix inequalities use \path{audit_identity.py}.
The complete list of $69$ triple types, the numerical inequalities in Lemmas~\ref{lem:sign} and~\ref{lem:class},
and the constants in Lemma~\ref{lem:rigidity} are checked by \path{audit_geometry.py}.
Proposition~\ref{prop:local} uses \path{audit_local.py}, the Taylor algebra in \path{localmin_jet.py},
and the rational proposals in \path{Certificates/Local/local_norm_bounds.json}.
The inner-node estimates and local ellipsoid containment tests use \path{prep_local3.py}.
Finally, \path{audit_closure.py} checks the expected result list, regenerated root data, problem parameters,
local/global overlap, tube radii and final rational comparisons. The analytic arguments connecting these
checks are given in the preceding sections.

\paragraph{Reproduction.}
From the repository root, install the pinned Python dependencies and run
\begin{verbatim}
python3 -m venv .venv
.venv/bin/python -m pip install -r requirements.txt
PYTHON="$PWD/.venv/bin/python" bash verify.sh 4
\end{verbatim}
The reference environment is Linux x86\_64 with Python 3.13.5, python-flint 0.9.0,
and g++ 14.2.0 with OpenMP. Bash and the standard utilities \texttt{sha256sum}, \texttt{head},
\texttt{tail}, \texttt{tee}, \texttt{grep} and \texttt{basename} are required; the Python standard
library reads the XZ archives. A TeX installation is needed only to rebuild this paper.
The README gives the full commands and the expected results.

The compilation flags are \texttt{-O2}, \texttt{-fopenmp}, \texttt{-ffp-contract=off} and \texttt{-fno-fast-math}.
The script uses the requested $1$--$8$ OpenMP threads
and one BLAS thread, and limits each process to $8\,000\,000$ KiB of virtual memory.
It writes generated inputs, binaries and new logs to \path{build/}; the supplied reference logs remain
in \path{results/reference/}. Python assertions remain enabled.
It stops on any failed command or assertion. The steps are:
\begin{enumerate}
\item verify the release hashes and the height, node-list, rigidity and local-minimality checks;
\item verify the exact matrix certificates and all energy identities;
\item regenerate the global root coverings and all $22$ local node calculations from the exact JSON data;
\item run all $31$ branch-and-bound calculations;
\item compare regenerated root data with the supplied archives and check every final coverage and margin inequality.
\end{enumerate}
All $31$ runs verify, totalling $2\,903\,200$ cells, with no failed cells or reported counterexamples.
The final success message is \texttt{FULL PROOF VERIFIED}. The corresponding marker is
\path{build/logs/FULL_PROOF_VERIFIED}, written only after the closure check succeeds.
The rational local coordinate systems and norm proposals are part of the supplied data, and every required
inequality is checked. Reproduction does not call an eigensolver. The reference run uses a fresh virtual environment
containing only python-flint. A root-data mismatch is a failed reproduction and must be investigated;
it must not be ignored to obtain a success message.

\paragraph{Scope of the computation.}
All inequalities used in the proof are established by exact rational tests or outward interval bounds.
The trusted software consists of the rational arithmetic routines, FLINT/Arb, the stated interval verifiers, and their
compiler and arithmetic environment, including correctly rounded binary64 operations and square root.
The proof is an analytic argument supplemented by these computations, rather than a formal proof-assistant derivation.
Neither convergence of a numerical optimizer nor agreement at finitely many sample configurations is used to establish
an identity or a positivity statement. The repository contains the fixed certificates and complete verification
sources; optimization and search software is not needed for reproduction.

\begin{thebibliography}{9}
\bibitem{CohnWoo} H.~Cohn, J.~Woo, \emph{Three-point bounds for energy minimization}, J. Amer. Math. Soc. 25 (2012), 929--958. \href{https://arxiv.org/abs/1103.0485}{arXiv:1103.0485}.
\bibitem{BV} C.~Bachoc, F.~Vallentin, \emph{New upper bounds for kissing numbers from semidefinite programming},
J. Amer. Math. Soc. 21 (2008), 909--924. \href{https://arxiv.org/abs/math/0608426}{arXiv:math/0608426}.
\end{thebibliography}
\end{document}
